\documentclass [10pt]{article}
\usepackage[margin=2cm]{geometry}
\usepackage[colorlinks=true,linkcolor=blue,urlcolor=blue]{hyperref}
\usepackage{listings}
\usepackage{graphicx}
\usepackage{tabulary}
\usepackage{longtable}
\usepackage{multirow}
\usepackage{hyperref}
\usepackage{mathtools}
\usepackage[all]{xy}

\title {Meat Language Reference Standard}
\author {Ron R Wills \textless\href{mailto:ron.rwsoft@gmail.com}{ron.rwsoft@gmail.com}\textgreater}

\begin {document}
\pagenumbering {gobble}
\begin{titlepage}
	\centering
	\includegraphics[width=0.35\textwidth]{meat-logo.jpg}\par\vspace{2cm}
	{\scshape\Large Language Version 1.0 Proposal\par}
	\vspace{1.5cm}
	{\huge\bfseries Meat Language Standard Reference\par}
	\vspace{2cm}
	{\Large\itshape Ron R Wills \texttt{\textless\href{mailto:ron.rwsoft@gmail.com}{ron.rwsoft@gmail.com}\textgreater}}\par

	\vfill

% Bottom of the page
	{\large \today\par}
\end{titlepage}

\noindent{}Copyright \copyright{} 2017 Ron R Wills \texttt{\textless\href{mailto:ron.rwsoft@gmail.com}{ron.rwsoft@gmail.com}\textgreater}.\\[.5 cm]
Permission is granted to copy, distribute and/or modify this document
under the terms of the GNU Free Documentation License, Version 1.3
or any later version published by the Free Software Foundation;
with no Invariant Sections, no Front-Cover Texts, and no Back-Cover Texts.
A copy of the license is included in the section entitled \textit{"GNU
Free Documentation License"}.
\tableofcontents
\newpage
\pagenumbering {arabic}

\section {Introduction}
\subsection {Scope}
The purpose of this document is not to be a learning guide to the language but
to define the Meat language itself.

\newpage
\section {Core Language Components}

All the components of the language described in this section are considered
mandatory to the language.

\subsection {Byte Interpretation}

\subsubsection{Integer Values}

\begin{center}
  \begin{tabular}{ r | c r r }
    & Size (bits) & Original Value & Endian Swapped Value\\ \hline
    Byte & 8 & $12_{16}$ & $12_{16}$ \\
    Word & 16 & $1234_{16}$ & $3412_{16}$ \\
    Integer & 32 & $12345678_{16}$ & $78563412_{16}$ \\
    Long & 64 & $\mathrm{123456789ABCDEF0}_{16}$ &
      $\mathrm{F0DEBC9A78563412}_{16}$\\
  \end{tabular}
\end{center}

\subsection{Virtual Machine Bytecode}

\begin{center}
  \begin{tabular}{ | p{6cm} | c | p{8cm} | }
    \hline
    Operation & Code & Brief Description \\ \hline
    No Operation & $00_{16}$ & This operation does nothing. \\
    Message Object & $01_{16}$ & Send a message to an object. \\
    Message Object with Result & $02_{16}$ & Send a message to the object
    and assigning the result to the context object. \\
    Message Object's Super Class & $03_{16}$ & Send a message to the super
    class.\\
    Message Object's Super Class with Result & $04_{16}$ & ... \\
    Beginning of Block Context & $\mathrm{0A}_{16}$ & ... \\
    End of Context & $\mathrm{0B}_{16}$ & ... \\
    Assign Local Variable & $10_{16}$ & ...\\
    Assign Object Property & $11_{16}$ & ...\\
    Assign Class Property & $12_{16}$ & ...\\
    Assign Class to Local Variable & $13_{16}$ & ...\\
    Assign an Integer Constant to a Local Variable & $14_{16}$ & ...\\
    \hline
  \end{tabular}
\end{center}

\subsubsection{Code $00_{16}$ -- No Operation}
The virtual machine only increased the byte code pointer for this byte code.
It's potental use is to align or pad byte codes if required.

The code pointer is increamented by one.

\subsubsection{Code $01_{16}$ -- Message Object}

\begin{center}
  \begin{tabular}{|l|p{3cm}|p{3cm}|p{3cm}|p{3cm}|}
    \hline
    \textbf{byte} & \textbf{byte} & \textbf{integer} & \textbf{byte} &
    \textbf{array of bytes} \\
    \hline
    $01_{16}$ & local id of object to message & message hash id &
    number of parameters & an array of local ids for each parameter \\
    \hline
  \end{tabular}
\end{center}
The virtual machine takes the following steps when sending a message to the
object:
\begin{enumerate}
\item A context is created from the local object found at the local id and the
  message hash id.
\item From the number of parameter and the local id parameter array, the
  parameters are added to the context.
\item The code pointer is increamented by 7 plus the number of parameters and
  updated in the current context.
\item Control is passed to the new context object.
\end{enumerate}

\subsubsection{Code $02_{16}$ -- Message Object with Result}

\begin{center}
  \begin{tabular}{|l|p{2.7cm}|p{2.7cm}|p{2.7cm}|p{2.7cm}|p{2.7cm}|}
    \hline
    \textbf{byte} & \textbf{byte} & \textbf{byte} & \textbf{integer} &
    \textbf{byte} & \textbf{array of bytes} \\
    \hline
    $02_{16}$ & local id of object to message & local id destination for the
    result value & message hash id & number of parameters & an array of local
    ids for each parameter \\
    \hline
  \end{tabular}
\end{center}
The virtual machine takes the following steps when sending a message to the
object:
\begin{enumerate}
\item A context is created from the local object found at the local id and the
  message hash id.
\item From the number of parameter and the local id parameter array, the
  parameters are added to the context.
\item Set the local id destination for the result value in the context.
\item The code pointer is increamented by 8 plus the number of parameters and
  updated in the current context.
\item Control is passed to the new context object.
\end{enumerate}

\subsection{Library File Format}

\begin{center}
  \begin{tabular}{ |l|p{4.5cm}|p{7.5cm}| }
    \hline
    \multirow{5}{*}{Header} & \texttt{unsigned bytes [4]} = ``MLIB''
    = $\mathrm{4D4C4942}_{16}$ & The
    magic file identifier\\ \cline{2-3}
    & \texttt{unsigned byte} = 1 & Major Meat library format version\\
    \cline{2-3}
    & \texttt{unsigned byte} = 0 & Minor Meat library format version\\
    \cline{2-3}
    & \texttt{unsigned word} = 0 & Reserved for future flags.\\ \cline{2-3}
    & \texttt{unsigned integer} = 0 & A hash ID for an application class for
    the virtual interpreter to message if the library is being executed.
    If this field is set to zero, then the library is not considered an
    executable application, but can still be loaded as a library.

    The hash ID must refer to an Application class within the same library.
    Any other reference is considered invalid and an Exception must be
    thrown to indicate the error by the virtual interpreter.\\
    \hline

    \multirow{2}{*}{Library Imports} & \texttt{unsigned byte} & The total number of libraries that need to be imported.\\ \cline{2-3}
    & \texttt{unsigned byte arrays} & A list of ''0'' terminated strings of each library name to be imported.\\ \hline

    \multirow{2}{*}{Serialized Class Objects} & \texttt{unsigned byte} & The total number of Classes serialized in the library.\\ \cline{2-3}
    & \texttt{serialized class array} & An array of all the serialized classes. see \hyperref[sec:class_serial]{Class Serialization}.\\ \hline
  \end{tabular}
\end{center}

\subsubsection{Class Serialization}
\label{sec:class_serial}

\textbf{Serialized Class Format}
\begin{center}
  \begin{tabular}{|l|l|p{8cm}|}
    \hline
    \multirow{3}{*}{Header} & \texttt{unsigned integer} & The hashed ID of the class name.\\ \cline{2-3}
    & \texttt{unsigned byte} & The number of class properties.\\ \cline{2-3}
    & \texttt{unsigned byte} & The number of object properties.\\ \hline

    \multirow{2}{*}{Virtual Table} & \texttt{unsigned byte} & The number of virtual table entries.\\ \cline{2-3}
    & \texttt{unsigned byte} & The number of class virtual table entries.\\ \hline

    \multirow{2}{*}{Byte Code} & \texttt{unsigned word} & Size of the byte code array.\\ \cline{2-3}
    & \texttt{unsigned byte array} & The byte code.\\ \hline
  \end{tabular}
\end{center}

\textbf{Virtual Table Entries}
\begin{center}
  \begin{tabular}{|l|p{8cm}|}
    \hline
    \texttt{unsigned integer} & Hashed ID of the method name.\\ \hline
    \texttt{unsigned integer} & Hashed ID of the class where the method is found.\\ \hline
    \texttt{unsigned byte} & Entry flags.\\ \hline
    \texttt{unsigned byte} & The number of local variables required during execution.\\ \hline
    \texttt{unsigned word} & Byte code offset to the start of the method.\\ \hline
  \end{tabular}
\end{center}

\textbf{Virtual Table Entry Flags}
\begin{displaymath}
  \xymatrix @-1pc {Bit &*+[F]{7} \ar @{-} [ddd]&*+[F]{6} \ar @{-} [dd] &*+[F]{5} \ar @{-} [dd] &*+[F]{4} \ar @{-} [dd] &*+[F]{3} \ar @{-} [dd] &*+[F]{2} \ar @{-} [dd]
    &*+[F]{1} \ar @{-} [dd] &*+[F]{0} \ar @{-} [d] & \\
    & & & & & & & & \ar [r] & \txt<6cm>{If set the virtual table entry is an offset into the byte code, otherwise it is an internal function pointer.}\\
    & & \ar [rrrrrrr] & & & & & & & \txt<6cm>{Reserved for future use.}\\
    & \ar [rrrrrrrr] & & & & & & & & \txt<6cm>{Method found in a super class.}}
\end{displaymath}


\newpage
\section {Core Language Builtin Library}
\subsection {Object Class}

The Object class is the base class for all objects and classes in the
Meat virtual machine.

\begin{lstlisting}
Object subClass: Object body: {
  class method == other
  class method != other
  class method is: other
  class method isNot: other
  class method isObject
  class method isClass
  class method isNull
  class method type
  method initialize
  method == other
  method != other
  method is: other
  method isNot: other
  method isObject
  method isClass
  method isNull
  method type
  method isType: class
  method weak
}
\end{lstlisting}

\subsubsection {The == methods}
\begin{lstlisting}
  class method == other
  method == other
\end{lstlisting}

Tests if the \textit{other} object or class is the same object as self. For
the Object class, it does not test the equality of values, but sub-classes
should override this method to test for value equality. The result is to be
either the True or False Boolean objects.

\hfill
\subsubsection {The != methods}
\begin{lstlisting}
  class method != other
  method != other
\end{lstlisting}

Tests if the \textit{other} object or class is not the same object as self.
For this class, it does not test the inequality of values, but sub-classes
should override this method to test for value inequality. The result is to be
either the True or False Boolean objects.

\hfill
\subsubsection {The is: methods}
\begin{lstlisting}
  class method is: other
  method is: other
\end{lstlisting}

Tests if the \textit{other} object or class is the same object as self. This
method should not be overridden by subclasses. The result is to be either the
True or False Boolean objects.

\hfill
\subsubsection {The isNot: methods}
\begin{lstlisting}
  class method isNot: other
  method isNot: other
\end{lstlisting}

Tests if the \textit{other} object or class is not the same object as self.
This method should not be overridden by subclasses. The result is to be either
the True or False Boolean objects.

\hfill
\subsubsection {The isObject methods}
\begin{lstlisting}
  class method isObject
  method isObject
\end{lstlisting}

Tests if the object is an object. The result is to be either the True or False
Boolean objects.

The class method always returns the False object and the object method always
returns the True object.

\hfill
\subsubsection {The isClass methods}
\begin{lstlisting}
  class method isClass
  method isClass
\end{lstlisting}

Tests if the object is a class. The result is to be either the True or False
Boolean objects.

The class method always returns the True object and the object method always
returns the False object.

\hfill
\subsubsection {The isNull methods}
\begin{lstlisting}
  class method isNull
  method isNull
\end{lstlisting}

Tests if the object or class is the Null object. These methods always return
the False object and should only be overridden by the Null class.

\hfill
\subsubsection {The type methods}
\begin{lstlisting}
  class method type
  method type
\end{lstlisting}

Returns the class for the object. The class method always returns the class
Class.

\hfill
\subsubsection {method initialize}
When a class creates a new object (see newObject) this method is messaged
to populate the objects properties with sane default values. This method
should be overridden by every sub-class and must always message the super
class initialize method.

\begin{lstlisting}
  method initialize {
    super initialize
    # Populate object's properties here.
  }
\end{lstlisting}

\subsubsection {method weak}

Returns a weak reference to the object that does not affect the normal memory
management reference counting.

\newpage
\subsection{Class Class}

The class Class is the base type for all Class objects.

\begin{lstlisting}
Object subClass: Class body: {
  class method subClass: name body: definition
  class method newObject
  class method superClass
  class method initialize
  method getName
}
\end{lstlisting}

\subsubsection{class method subClass: name body: definition}
This method should raise an Exception. The functionality for this method
is to be provided by a compiler library. How a compiler library hooks in the
functionaliy is up to the underlying implementation and is not defined here.

\subsubsection{class method newObject}
\label{sec:class_method_newobject}
Method used for creating constructor class methods. This method creates a new
object and sets the self object to it then messages \textit{self initialize}.

All constructor methods should follow the following template:

\begin{lstlisting}
class method new {
  self newObject
  # Add any additional object initialization here.
  context return: self
}
\end{lstlisting}

\subsubsection{class method superClass}
Returns the super class for the Class.

\subsubsection{class method initialize}
When a library is loaded and the class is unserialized, this method is
messaged to populate any class properties with default values. Any class
properties that are not set by this method will default to the Null object.
This method can be overridden by sub-classes and must always message the
super class initialize method.

\begin{lstlisting}
  method initialize {
    super initialize
    # Populate object's properties here.
  }
\end{lstlisting}

\subsubsection{method getName}
Returns the Text representation of the class name.

\subsection {Context Class}

A context is the execution context when a method answers a message.
\begin{center}
  \begin{tabular}{ c | l }
    Local Variable & Value\\ \hline
    0 & self \\
    1 & class \\
    2 & context (weak reference) \\
    3 & null \\
    4 \dots & message parameters \\
    \dots last & local variables \\
  \end{tabular}
\end{center}

\begin{lstlisting}
Object subClass: Context body: {
  class method subClass: name body: definition
  class method newObject
  method messenger
  method localVariables
  method getLocal: index
  method setLocal: index to: value
  method return
  method return: result
}
\end{lstlisting}

\subsubsection {class method subClass: name body: definition}
The Context class should never be sub-classed therefore this method always
raises an exception.

\hfill
\subsubsection {class method newObject}
Context classes are only created internally when an object is messaged
therefore this method always raises an exception.

\hfill
\subsubsection {method messenger}
This method returns the context object where the message was
initiated. The null object can be returned if it was the very first context
created to start execution.

\subsection {BlockContext Class}

\begin{displaymath}
  \xymatrix @-1pc {
    & \txt{Origin Context} \ar [dl] \ar@{<=>}[dd] \\
    \txt{Messaging Context} \ar@{<->}[dr] & \\
    & \txt{Block Context}\\
  }
\end{displaymath}

The origin context is the context in which the block context byte code was found and the block context was created. The local variables for the origin context are prepended to the block context local variables.

The messaging context is the context where one of the \textit{execute} method is messaged.

\begin{center}
  \begin{tabular}{ c | l }
  	Variable   & Context        \\ \hline
  	self       & Origin Context \\
  	class      & Origin Context \\
  	context    & Block Context  \\
  	null       & Origin Context \\
  	4 \dots    & Origin Context \\
  	\dots last & Block Context
  \end{tabular}
\end{center}


\begin{lstlisting}
Context subClass: BlockContext body: {
  method execute
  method executeOnBreak: block
  method executeOnBreak: breakBlock onContinue: continueBlock
  method executeOnContinue: block
  method localVariables
  method getLocal: index
  method setLocal: index to: value
  method return
  method return: result
  method continue
  method break
}
\end{lstlisting}

\subsubsection{method execute}
Executes the block. After execution the block must be reset to
allow the block to be executed again.

\subsubsection{method executeOnBreak: block}

\subsubsection{method executeOnBreak: breakBlock onContinue: continueBlock}

\subsubsection{method executeOnContinue: block}

\subsubsection{method localVariables}

\subsubsection{method getLocal: index}

\subsubsection{method setLocal: index to: value}

\subsubsection{method return}
Stops execution of this context and messages the \textit{return} method of the origin context.

\subsubsection{method return: result}
Stops execution of this context and messages the \textit{return: value} method of the origin context. The \textit{result} is not set to the block context.

\subsubsection{continue}

\subsubsection{break}
If the Block Context was messaged with \textit{executeOnBreak:} or \textit{executeOnBreak:onContinue:} then the Block Context ends execution and returns to the Messaging Context to execute appropriate break block.

If the Block Context was messaged with \textit{execute} or \textit{executeonContinue:} then \textit{break} is messaged to the origin context.

\subsection {Null Class and Object}

The Null object is to indicate that a reference doesn't actually reference
any object. The Null class has only one object to define null and can not
be subclassed.

\begin{lstlisting}
Object subclass: Null body: {
  class method newObject
  class method subClass: name body: definition
  class method object
  method isNull
}
\end{lstlisting}

\subsubsection{Class method newObject}
Always raises an Exception.

\subsubsection{Class method subClass: body:}
Always raises an Exception.

\subsubsection{Class method object}
The method returns the null object.

\subsubsection{Method isNull}
The method returns the true object.

\subsection {Exception Class}

\begin{lstlisting}
Object subclass: Exception body: {
  class method try: block
  class method try: block catch: catchBlock
  class method try: block catch: .error do: catchBlock
  class method throw
  class method throw: message
  class method throwFor: context
  class method throw: message for: context
  method message
  method context
}
\end{lstlisting}

\subsubsection{Class method try:}

\subsection {Boolean Class and Objects}

\begin{lstlisting}
Object subClass: Boolean body: {
  class method true
  class method false
  method isTrue: block
  method isTrue: trueBlock else: falseBlock
  method isFalse: block
  method isFalse: falseBlock else: trueBlock
  method not
  method and: other
  method or: other
  method xor: other
  method asText
}
\end{lstlisting}

\subsection{Numeric Class}

The Numeric Class is an abstract base class for all numeric or mathematical values.

\begin{lstlisting}
Object subClass: Numeric body: {
  class method minValue
  class method maxValue
  method + other
  method - other
  method * other
  method / other
  method % other
  method ^ other
  method == other
  method <> other
  method < other
  method <= other
  method > other
  method >= other
  method negative
  method absolute
  method sign
  method asText
}
\end{lstlisting}

\subsubsection{class method minValue}
This method is meant to be overridden by sub-classes to return a object of the smallest value possible for the numeric type. For this class this method just raises an Exception.

\subsubsection{class method maxValue}
This method is meant to be overridden by sub-classes to return a object of the largest value possible for the numeric type. For this class this method just raises an Exception.

\subsubsection{method + other}

\subsubsection{method - other}

\subsubsection{method * other}

\subsubsection{method / other}

\subsubsection{method \% other}
	
\subsubsection{method \textasciicircum other}

\subsubsection{method $==$ other}

\subsubsection{method $<>$ other}

\subsubsection{method $<$ other}

\subsubsection{method $<=$ other}

\subsubsection{method $>$ other}

\subsubsection{method $>=$ other}

\subsubsection{negative}

\subsubsection{absolute}

\subsubsection{sign}

\subsection{Integer Class}

\begin{lstlisting}
Numeric subClass: Integer body: {
  class method minValue
  class method maxValue
  method + other
  method - other
  method * other
  method / other
  method % other
  method ^ other
  method == other
  method != other
  method < other
  method <= other
  method > other
  method >= other
  method negative
  method absolute
  method sign
  method or: other
  method and: other
  method xor: other
  method rshift
  method rshift: amount
  method lshift
  method lshift: amount
  mothod timesDo: block
  method asText
}
\end{lstlisting}

\subsection{Number Class}

\begin{lstlisting}
Numeric subclass: Number body: {
  class method cos: angle
  
}
\end{lstlisting}
\subsection {Text Class}

\begin{lstlisting}
Object subClass: Text body: {
  method isEmpty
  method length
  method getCharAt: index
  method getTextFrom: start to: end
  method changeFrom: start to: end with: value
  method replaceAll: text with: value
  method findFirst: other
  method findFirst: other at: index
  method findLast: other
  method forEach: character do: block
  method asLowercase
  method asUppercase
  method copy
  method asNumber
  method asInteger
  method sameAs: other
  method == other
  method != other
  method < other
  method > other
  method + other
  method * amount
}
\end{lstlisting}

\subsubsection{method isEmpty}

\subsubsection{method length}
Returns an Integer object with the number of unicode characters in the text.

\subsection{List Class}

\begin{lstlisting}
Object subClass: List body: {
  method copy
}
\end{lstlisting}

\subsubsection{method copy}

\subsection {Set Class}

\begin{lstlisting}
Object subclass: Set body: {
  class method new
}
\end{lstlisting}
\subsection {Index Class}

\begin{lstlisting}
Object subClass: Index body: {
  class method new
  method get: key
  method set: key to: value
  method has: key
  method remove: key
  method entries
  method getKeyList
  method getValueList
  method clear
  method isEmpty
  method copy
  method forEach: key with: value do: body
  method forEachKey: key do:
  method forEachValue: value do:
}
\end{lstlisting}

\subsection{Archive Class}

\begin{lstlisting}
Object subclass: Archive body: {

}
\end{lstlisting}
\subsection{Library Class}

\begin{lstlisting}
Object subClass: Library body: {
  class method exists: libraryName
  class method import: libraryName
  class method requires: libraryName
  class method include: nativeCode
  class method setApplication: class
  class method unload: libraryName
  class method listSearchPaths
}
\end{lstlisting}

\subsubsection{class method exists:}
Checks every search path directory for the library given by \textit{libraryName} and return the true Boolean object if the library was found otherwise it returns the false Boolean object.

\subsubsection{class method import:}

\subsubsection{class method requires:}
By default this method always throws an Exception. This method should be hooked when a compiler is loaded.
\subsection{Application Class}

\newpage
\section {Grinder Lexical Elements and Parsing}
\subsection {Characters}
The Unicode UTF-8 character set

The Unicode white space characters \texttt{U+0009}
\textit{(the TAB character \textbackslash t)}, \texttt{U+0020}
\textit{(the SPACE character)} and
\texttt{U+00A0} \textit{(the NO-BREAK SPACE character)} are used as delimiters
between the token components of a command. These characters can be used
multiple times together to form a single delimiter.

A line termination or a \textit{newline} is definied by the character set
described in the following table. Code parsers must accept all of these
characters as a line terminator.

\begin{center}
  \begin{tabular}{|l|l|l|}
  	\hline
  	                          & \textbf{Unicode}                & \textbf{Escaped Character}        \\ \hline
  	\textbf{LF}               & \texttt{U+000A}                 & \textbackslash n                  \\ \hline
  	\textbf{VT}               & \texttt{U+000B}                 & \textbackslash v                  \\ \hline
  	\textbf{FF}               & \texttt{U+000C}                 & \textbackslash f                  \\ \hline
  	\textbf{CR}               & \texttt{U+000D}                 & \textbackslash r                  \\ \hline
  	\textbf{CR} + \textbf{LF} & \texttt{U+000D} \texttt{U+000A} & \textbackslash r \textbackslash n \\ \hline
  	\textbf{LS}               & \texttt{U+2028}                 &  \\ \hline
  	\textbf{PS}               & \texttt{U+2029}                 &  \\ \hline
  \end{tabular}
\end{center}

\subsubsection {Numeric Chararcters}

Using the unicode character set allows code to be in any language supported.
This also extends to numeric character sets. The following table describes the
numeric characters that meat parsers must recognize as numeric values and
translate into their binary equivalent.

\begin{center}
  \begin{longtable}{|l|c|}
  	\hline
  	                                             &         \textbf{Unicode}          \\ \hline
  	\textbf{Digits}                              &  \texttt{U+0030}-\texttt{U+0039}  \\ \hline
  	\textbf{Arabic-Indic Digits}                 &  \texttt{U+0660}-\texttt{U+0669}  \\ \hline
  	\textbf{Extended Arabic-Indic Digit}         &  \texttt{U+06F0}-\texttt{U+06F9}  \\ \hline
  	\textbf{NKo Digits}                          &  \texttt{U+07C0}-\texttt{U+07C9}  \\ \hline
  	\textbf{Devavagari Digits}                   &  \texttt{U+0966}-\texttt{U+096F}  \\ \hline
  	\textbf{Bengali Digits}                      &  \texttt{U+09E6}-\texttt{U+09EF}  \\ \hline
  	\textbf{Gurmukhi Digits}                     &  \texttt{U+0A66}-\texttt{U+0A6F}  \\ \hline
  	\textbf{Gujarati Digits}                     &  \texttt{U+0AE6}-\texttt{U+0AEF}  \\ \hline
  	\textbf{Oriya Digits}                        &  \texttt{U+0B66}-\texttt{U+0B6F}  \\ \hline
  	\textbf{Tamil Digits}                        &  \texttt{U+0BE6}-\texttt{U+0BEF}  \\ \hline
  	\textbf{Telugu Digits}                       &  \texttt{U+0C66}-\texttt{U+0C6F}  \\ \hline
  	\textbf{Kannada Digits}                      &  \texttt{U+0CE6}-\texttt{U+0CEF}  \\ \hline
  	\textbf{Malayalam Digits}                    &  \texttt{U+0D66}-\texttt{U+0D6F}  \\ \hline
  	\textbf{Sinhala Digits}                      &  \texttt{U+0DE6}-\texttt{U+0DEF}  \\ \hline
  	\textbf{Thai Digits}                         &  \texttt{U+0E50}-\texttt{U+0E59}  \\ \hline
  	\textbf{Lao Digits}                          &  \texttt{U+0ED0}-\texttt{U+0ED9}  \\ \hline
  	\textbf{Tibetan Digits}                      &  \texttt{U+0F20}-\texttt{U+0F29}  \\ \hline
  	\textbf{Myanmar Digits}                      &  \texttt{U+1040}-\texttt{U+1049}  \\ \hline
  	\textbf{Myanmar Shan Digits}                 &  \texttt{U+1090}-\texttt{U+1099}  \\ \hline
  	\textbf{Khmer Digits}                        &  \texttt{U+17E0}-\texttt{U+17E9}  \\ \hline
  	\textbf{Mongolian Digits}                    &  \texttt{U+1810}-\texttt{U+1819}  \\ \hline
  	\textbf{Limbu Digits}                        &  \texttt{U+1946}-\texttt{U+194F}  \\ \hline
  	\textbf{New Tai Lue Digits}                  &  \texttt{U+19D0}-\texttt{U+19D9}  \\ \hline
  	\textbf{Tai Tham Hora Digits}                &  \texttt{U+1A80}-\texttt{U+1A89}  \\ \hline
  	\textbf{Tai Tham Tham Digits}                &  \texttt{U+1A90}-\texttt{U+1A99}  \\ \hline
  	\textbf{Balinese Digits}                     &  \texttt{U+1B50}-\texttt{U+1B59}  \\ \hline
  	\textbf{Sundanese Digits}                    &  \texttt{U+1BB0}-\texttt{U+1BB9}  \\ \hline
  	\textbf{Lepcha Digits}                       &  \texttt{U+1C40}-\texttt{U+1C49}  \\ \hline
  	\textbf{Ol Chiki Digits}                     &  \texttt{U+1C50}-\texttt{U+1C59}  \\ \hline
  	\textbf{Vai Digits}                          &  \texttt{U+A620}-\texttt{U+A629}  \\ \hline
  	\textbf{Saurashtra Digits}                   &  \texttt{U+A8D0}-\texttt{U+A8D9}  \\ \hline
  	\textbf{Kayah Li Digits}                     &  \texttt{U+A900}-\texttt{U+A909}  \\ \hline
  	\textbf{Javanese Digits}                     &  \texttt{U+A9D0}-\texttt{U+A9D9}  \\ \hline
  	\textbf{Myanmar Tai Laing Digits}            &  \texttt{U+A9F0}-\texttt{U+A9F9}  \\ \hline
  	\textbf{Cham Digits}                         &  \texttt{U+AA50}-\texttt{U+AA59}  \\ \hline
  	\textbf{Meetei Mayek Digits}                 &  \texttt{U+ABF0}-\texttt{U+ABF9}  \\ \hline
  	\textbf{Fullwidth Digits}                    &  \texttt{U+FF10}-\texttt{U+FF19}  \\ \hline
  	\textbf{Osmanya Digits}                      & \texttt{U+104A0}-\texttt{U+104A9} \\ \hline
  	\textbf{Brahmi Digits}                       & \texttt{U+11066}-\texttt{U+1106F} \\ \hline
  	\textbf{Sora Sompeng Digits}                 & \texttt{U+110F0}-\texttt{U+110F9} \\ \hline
  	\textbf{Chakma Digits}                       & \texttt{U+11136}-\texttt{U+1113F} \\ \hline
  	\textbf{Sharada Digits}                      & \texttt{U+111D0}-\texttt{U+111D9} \\ \hline
  	\textbf{Khudawadi Digits}                    & \texttt{U+112F0}-\texttt{U+112F9} \\ \hline
  	\textbf{Newa Digits}                         & \texttt{U+11450}-\texttt{U+11459} \\ \hline
  	\textbf{Tirhuta Digits}                      & \texttt{U+114D0}-\texttt{U+114D9} \\ \hline
  	\textbf{Modi Digits}                         & \texttt{U+11650}-\texttt{U+11659} \\ \hline
  	\textbf{Takri Digits}                        & \texttt{U+116C0}-\texttt{U+116C9} \\ \hline
  	\textbf{Ahom Digits}                         & \texttt{U+11730}-\texttt{U+11739} \\ \hline
  	\textbf{Warang Citi Digits}                  & \texttt{U+118E0}-\texttt{U+118E9} \\ \hline
  	\textbf{Bhaiksuki Digits}                    & \texttt{U+11C50}-\texttt{U+11C59} \\ \hline
  	\textbf{Masaram Digits}                      & \texttt{U+11D50}-\texttt{U+11D59} \\ \hline
  	\textbf{Mro Digits}                          & \texttt{U+16A60}-\texttt{U+16A69} \\ \hline
  	\textbf{Pahawh Hmong Digits}                 & \texttt{U+16B50}-\texttt{U+16B59} \\ \hline
  	\textbf{Mathematical Bold Digits}            & \texttt{U+1D7CE}-\texttt{U+1D7D7} \\ \hline
  	\textbf{Mathematical Double-Struck Digits}   & \texttt{U+1D7D8}-\texttt{U+1D7EF} \\ \hline
  	\textbf{Mathematical Sans-Serif Digits}      & \texttt{U+1D7E2}-\texttt{U+1D7EB} \\ \hline
  	\textbf{Mathematical Sans-Serif Bold Digits} & \texttt{U+1D7EC}-\texttt{U+1D7F5} \\ \hline
  	\textbf{Mathematical Monospace Digits}       & \texttt{U+1D7F6}-\texttt{U+1D7FF} \\ \hline
  	\textbf{Adlam Digits}                        & \texttt{U+1E950}-\texttt{U+1E959} \\ \hline
  \end{longtable}
\end{center}

\newpage
\section {Appendix}
\subsection {GNU Free Documentation License}
\begin{center}
  Version 1.3, 3 November 2008\\
  Copyright \copyright{} 2000, 2001, 2002, 2007, 2008  Free Software Foundation, Inc.\\
 \bigskip
 \texttt{<http://fsf.org/>}\\
 \bigskip
\end{center}

\noindent{}Everyone is permitted to copy and distribute verbatim copies
of this license document, but changing it is not allowed.

\begin{center}
{\bf\large Preamble}
\end{center}

\noindent{}The purpose of this License is to make a manual, textbook, or other
functional and useful document ``free'' in the sense of freedom: to
assure everyone the effective freedom to copy and redistribute it,
with or without modifying it, either commercially or noncommercially.
Secondarily, this License preserves for the author and publisher a way
to get credit for their work, while not being considered responsible
for modifications made by others.

This License is a kind of ``copyleft'', which means that derivative
works of the document must themselves be free in the same sense.  It
complements the GNU General Public License, which is a copyleft
license designed for free software.

We have designed this License in order to use it for manuals for free
software, because free software needs free documentation: a free
program should come with manuals providing the same freedoms that the
software does.  But this License is not limited to software manuals;
it can be used for any textual work, regardless of subject matter or
whether it is published as a printed book.  We recommend this License
principally for works whose purpose is instruction or reference.

\subsubsection {APPLICABILITY AND DEFINITIONS}

This License applies to any manual or other work, in any medium, that
contains a notice placed by the copyright holder saying it can be
distributed under the terms of this License.  Such a notice grants a
world-wide, royalty-free license, unlimited in duration, to use that
work under the conditions stated herein.  The ``\textbf{Document}'', below,
refers to any such manual or work.  Any member of the public is a
licensee, and is addressed as ``\textbf{you}''.  You accept the license if you
copy, modify or distribute the work in a way requiring permission
under copyright law.

A ``\textbf{Modified Version}'' of the Document means any work containing the
Document or a portion of it, either copied verbatim, or with
modifications and/or translated into another language.

A ``\textbf{Secondary Section}'' is a named appendix or a front-matter section of
the Document that deals exclusively with the relationship of the
publishers or authors of the Document to the Document's overall subject
(or to related matters) and contains nothing that could fall directly
within that overall subject.  (Thus, if the Document is in part a
textbook of mathematics, a Secondary Section may not explain any
mathematics.)  The relationship could be a matter of historical
connection with the subject or with related matters, or of legal,
commercial, philosophical, ethical or political position regarding
them.

The ``\textbf{Invariant Sections}'' are certain Secondary Sections whose titles
are designated, as being those of Invariant Sections, in the notice
that says that the Document is released under this License.  If a
section does not fit the above definition of Secondary then it is not
allowed to be designated as Invariant.  The Document may contain zero
Invariant Sections.  If the Document does not identify any Invariant
Sections then there are none.

The ``\textbf{Cover Texts}'' are certain short passages of text that are listed,
as Front-Cover Texts or Back-Cover Texts, in the notice that says that
the Document is released under this License.  A Front-Cover Text may
be at most 5 words, and a Back-Cover Text may be at most 25 words.

A ``\textbf{Transparent}'' copy of the Document means a machine-readable copy,
represented in a format whose specification is available to the
general public, that is suitable for revising the document
straightforwardly with generic text editors or (for images composed of
pixels) generic paint programs or (for drawings) some widely available
drawing editor, and that is suitable for input to text formatters or
for automatic translation to a variety of formats suitable for input
to text formatters.  A copy made in an otherwise Transparent file
format whose markup, or absence of markup, has been arranged to thwart
or discourage subsequent modification by readers is not Transparent.
An image format is not Transparent if used for any substantial amount
of text.  A copy that is not ``Transparent'' is called ``\textbf{Opaque}''.

Examples of suitable formats for Transparent copies include plain
ASCII without markup, Texinfo input format, LaTeX input format, SGML
or XML using a publicly available DTD, and standard-conforming simple
HTML, PostScript or PDF designed for human modification.  Examples of
transparent image formats include PNG, XCF and JPG.  Opaque formats
include proprietary formats that can be read and edited only by
proprietary word processors, SGML or XML for which the DTD and/or
processing tools are not generally available, and the
machine-generated HTML, PostScript or PDF produced by some word
processors for output purposes only.

The ``\textbf{Title Page}'' means, for a printed book, the title page itself,
plus such following pages as are needed to hold, legibly, the material
this License requires to appear in the title page.  For works in
formats which do not have any title page as such, ``Title Page'' means
the text near the most prominent appearance of the work's title,
preceding the beginning of the body of the text.

The ``\textbf{publisher}'' means any person or entity that distributes
copies of the Document to the public.

A section ``\textbf{Entitled XYZ}'' means a named subunit of the Document whose
title either is precisely XYZ or contains XYZ in parentheses following
text that translates XYZ in another language.  (Here XYZ stands for a
specific section name mentioned below, such as ``\textbf{Acknowledgements}'',
``\textbf{Dedications}'', ``\textbf{Endorsements}'', or ``\textbf{History}''.)
To ``\textbf{Preserve the Title}''
of such a section when you modify the Document means that it remains a
section ``Entitled XYZ'' according to this definition.

The Document may include Warranty Disclaimers next to the notice which
states that this License applies to the Document.  These Warranty
Disclaimers are considered to be included by reference in this
License, but only as regards disclaiming warranties: any other
implication that these Warranty Disclaimers may have is void and has
no effect on the meaning of this License.

\subsubsection{VERBATIM COPYING}

You may copy and distribute the Document in any medium, either
commercially or noncommercially, provided that this License, the
copyright notices, and the license notice saying this License applies
to the Document are reproduced in all copies, and that you add no other
conditions whatsoever to those of this License.  You may not use
technical measures to obstruct or control the reading or further
copying of the copies you make or distribute.  However, you may accept
compensation in exchange for copies.  If you distribute a large enough
number of copies you must also follow the conditions in section
\arabic{section}.\arabic{subsection}.3.

You may also lend copies, under the same conditions stated above, and
you may publicly display copies.

\subsubsection{COPYING IN QUANTITY}


If you publish printed copies (or copies in media that commonly have
printed covers) of the Document, numbering more than 100, and the
Document's license notice requires Cover Texts, you must enclose the
copies in covers that carry, clearly and legibly, all these Cover
Texts: Front-Cover Texts on the front cover, and Back-Cover Texts on
the back cover.  Both covers must also clearly and legibly identify
you as the publisher of these copies.  The front cover must present
the full title with all words of the title equally prominent and
visible.  You may add other material on the covers in addition.
Copying with changes limited to the covers, as long as they preserve
the title of the Document and satisfy these conditions, can be treated
as verbatim copying in other respects.

If the required texts for either cover are too voluminous to fit
legibly, you should put the first ones listed (as many as fit
reasonably) on the actual cover, and continue the rest onto adjacent
pages.

If you publish or distribute Opaque copies of the Document numbering
more than 100, you must either include a machine-readable Transparent
copy along with each Opaque copy, or state in or with each Opaque copy
a computer-network location from which the general network-using
public has access to download using public-standard network protocols
a complete Transparent copy of the Document, free of added material.
If you use the latter option, you must take reasonably prudent steps,
when you begin distribution of Opaque copies in quantity, to ensure
that this Transparent copy will remain thus accessible at the stated
location until at least one year after the last time you distribute an
Opaque copy (directly or through your agents or retailers) of that
edition to the public.

It is requested, but not required, that you contact the authors of the
Document well before redistributing any large number of copies, to give
them a chance to provide you with an updated version of the Document.

\subsubsection{MODIFICATIONS}

You may copy and distribute a Modified Version of the Document under
the conditions of sections \arabic{section}.\arabic{subsection}.2 and \arabic{section}.\arabic{subsection}.3 above, provided that you release
the Modified Version under precisely this License, with the Modified
Version filling the role of the Document, thus licensing distribution
and modification of the Modified Version to whoever possesses a copy
of it.  In addition, you must do these things in the Modified Version:

\begin{itemize}
\item[A.]
   Use in the Title Page (and on the covers, if any) a title distinct
   from that of the Document, and from those of previous versions
   (which should, if there were any, be listed in the History section
   of the Document).  You may use the same title as a previous version
   if the original publisher of that version gives permission.

\item[B.]
   List on the Title Page, as authors, one or more persons or entities
   responsible for authorship of the modifications in the Modified
   Version, together with at least five of the principal authors of the
   Document (all of its principal authors, if it has fewer than five),
   unless they release you from this requirement.

\item[C.]
   State on the Title page the name of the publisher of the
   Modified Version, as the publisher.

\item[D.]
   Preserve all the copyright notices of the Document.

\item[E.]
   Add an appropriate copyright notice for your modifications
   adjacent to the other copyright notices.

\item[F.]
   Include, immediately after the copyright notices, a license notice
   giving the public permission to use the Modified Version under the
   terms of this License, in the form shown in the Addendum below.

\item[G.]
   Preserve in that license notice the full lists of Invariant Sections
   and required Cover Texts given in the Document's license notice.

\item[H.]
   Include an unaltered copy of this License.

\item[I.]
   Preserve the section Entitled ``History'', Preserve its Title, and add
   to it an item stating at least the title, year, new authors, and
   publisher of the Modified Version as given on the Title Page.  If
   there is no section Entitled ``History'' in the Document, create one
   stating the title, year, authors, and publisher of the Document as
   given on its Title Page, then add an item describing the Modified
   Version as stated in the previous sentence.

\item[J.]
   Preserve the network location, if any, given in the Document for
   public access to a Transparent copy of the Document, and likewise
   the network locations given in the Document for previous versions
   it was based on.  These may be placed in the ``History'' section.
   You may omit a network location for a work that was published at
   least four years before the Document itself, or if the original
   publisher of the version it refers to gives permission.

\item[K.]
   For any section Entitled ``Acknowledgements'' or ``Dedications'',
   Preserve the Title of the section, and preserve in the section all
   the substance and tone of each of the contributor acknowledgements
   and/or dedications given therein.

\item[L.]
   Preserve all the Invariant Sections of the Document,
   unaltered in their text and in their titles.  Section numbers
   or the equivalent are not considered part of the section titles.

\item[M.]
   Delete any section Entitled ``Endorsements''.  Such a section
   may not be included in the Modified Version.

\item[N.]
   Do not retitle any existing section to be Entitled ``Endorsements''
   or to conflict in title with any Invariant Section.

\item[O.]
   Preserve any Warranty Disclaimers.
\end{itemize}

If the Modified Version includes new front-matter sections or
appendices that qualify as Secondary Sections and contain no material
copied from the Document, you may at your option designate some or all
of these sections as invariant.  To do this, add their titles to the
list of Invariant Sections in the Modified Version's license notice.
These titles must be distinct from any other section titles.

You may add a section Entitled ``Endorsements'', provided it contains
nothing but endorsements of your Modified Version by various
parties---for example, statements of peer review or that the text has
been approved by an organization as the authoritative definition of a
standard.

You may add a passage of up to five words as a Front-Cover Text, and a
passage of up to 25 words as a Back-Cover Text, to the end of the list
of Cover Texts in the Modified Version.  Only one passage of
Front-Cover Text and one of Back-Cover Text may be added by (or
through arrangements made by) any one entity.  If the Document already
includes a cover text for the same cover, previously added by you or
by arrangement made by the same entity you are acting on behalf of,
you may not add another; but you may replace the old one, on explicit
permission from the previous publisher that added the old one.

The author(s) and publisher(s) of the Document do not by this License
give permission to use their names for publicity for or to assert or
imply endorsement of any Modified Version.

\subsubsection{COMBINING DOCUMENTS}

You may combine the Document with other documents released under this
License, under the terms defined in section \arabic{section}.\arabic{subsection}.4 above for modified
versions, provided that you include in the combination all of the
Invariant Sections of all of the original documents, unmodified, and
list them all as Invariant Sections of your combined work in its
license notice, and that you preserve all their Warranty Disclaimers.

The combined work need only contain one copy of this License, and
multiple identical Invariant Sections may be replaced with a single
copy.  If there are multiple Invariant Sections with the same name but
different contents, make the title of each such section unique by
adding at the end of it, in parentheses, the name of the original
author or publisher of that section if known, or else a unique number.
Make the same adjustment to the section titles in the list of
Invariant Sections in the license notice of the combined work.

In the combination, you must combine any sections Entitled ``History''
in the various original documents, forming one section Entitled
``History''; likewise combine any sections Entitled ``Acknowledgements'',
and any sections Entitled ``Dedications''.  You must delete all sections
Entitled ``Endorsements''.

\subsubsection{COLLECTIONS OF DOCUMENTS}

You may make a collection consisting of the Document and other documents
released under this License, and replace the individual copies of this
License in the various documents with a single copy that is included in
the collection, provided that you follow the rules of this License for
verbatim copying of each of the documents in all other respects.

You may extract a single document from such a collection, and distribute
it individually under this License, provided you insert a copy of this
License into the extracted document, and follow this License in all
other respects regarding verbatim copying of that document.

\subsubsection{AGGREGATION WITH INDEPENDENT WORKS}

A compilation of the Document or its derivatives with other separate
and independent documents or works, in or on a volume of a storage or
distribution medium, is called an ``aggregate'' if the copyright
resulting from the compilation is not used to limit the legal rights
of the compilation's users beyond what the individual works permit.
When the Document is included in an aggregate, this License does not
apply to the other works in the aggregate which are not themselves
derivative works of the Document.

If the Cover Text requirement of section \arabic{section}.\arabic{subsection}.3 is applicable to these
copies of the Document, then if the Document is less than one half of
the entire aggregate, the Document's Cover Texts may be placed on
covers that bracket the Document within the aggregate, or the
electronic equivalent of covers if the Document is in electronic form.
Otherwise they must appear on printed covers that bracket the whole
aggregate.

\subsubsection{TRANSLATION}

Translation is considered a kind of modification, so you may
distribute translations of the Document under the terms of section \arabic{section}.\arabic{subsection}.4.
Replacing Invariant Sections with translations requires special
permission from their copyright holders, but you may include
translations of some or all Invariant Sections in addition to the
original versions of these Invariant Sections.  You may include a
translation of this License, and all the license notices in the
Document, and any Warranty Disclaimers, provided that you also include
the original English version of this License and the original versions
of those notices and disclaimers.  In case of a disagreement between
the translation and the original version of this License or a notice
or disclaimer, the original version will prevail.

If a section in the Document is Entitled ``Acknowledgements'',
``Dedications'', or ``History'', the requirement (section \arabic{section}.\arabic{subsection}.4) to Preserve
its Title (section \arabic{section}.\arabic{subsection}.1) will typically require changing the actual
title.

\subsubsection{TERMINATION}

You may not copy, modify, sublicense, or distribute the Document
except as expressly provided under this License.  Any attempt
otherwise to copy, modify, sublicense, or distribute it is void, and
will automatically terminate your rights under this License.

However, if you cease all violation of this License, then your license
from a particular copyright holder is reinstated (a) provisionally,
unless and until the copyright holder explicitly and finally
terminates your license, and (b) permanently, if the copyright holder
fails to notify you of the violation by some reasonable means prior to
60 days after the cessation.

Moreover, your license from a particular copyright holder is
reinstated permanently if the copyright holder notifies you of the
violation by some reasonable means, this is the first time you have
received notice of violation of this License (for any work) from that
copyright holder, and you cure the violation prior to 30 days after
your receipt of the notice.

Termination of your rights under this section does not terminate the
licenses of parties who have received copies or rights from you under
this License.  If your rights have been terminated and not permanently
reinstated, receipt of a copy of some or all of the same material does
not give you any rights to use it.

\subsubsection{FUTURE REVISIONS OF THIS LICENSE}

The Free Software Foundation may publish new, revised versions
of the GNU Free Documentation License from time to time.  Such new
versions will be similar in spirit to the present version, but may
differ in detail to address new problems or concerns.  See
\texttt{http://www.gnu.org/copyleft/}.

Each version of the License is given a distinguishing version number.
If the Document specifies that a particular numbered version of this
License ``or any later version'' applies to it, you have the option of
following the terms and conditions either of that specified version or
of any later version that has been published (not as a draft) by the
Free Software Foundation.  If the Document does not specify a version
number of this License, you may choose any version ever published (not
as a draft) by the Free Software Foundation.  If the Document
specifies that a proxy can decide which future versions of this
License can be used, that proxy's public statement of acceptance of a
version permanently authorizes you to choose that version for the
Document.

\subsubsection{RELICENSING}

``Massive Multiauthor Collaboration Site'' (or ``MMC Site'') means any
World Wide Web server that publishes copyrightable works and also
provides prominent facilities for anybody to edit those works.  A
public wiki that anybody can edit is an example of such a server.  A
``Massive Multiauthor Collaboration'' (or ``MMC'') contained in the
site means any set of copyrightable works thus published on the MMC
site.

``CC-BY-SA'' means the Creative Commons Attribution-Share Alike 3.0
license published by Creative Commons Corporation, a not-for-profit
corporation with a principal place of business in San Francisco,
California, as well as future copyleft versions of that license
published by that same organization.

``Incorporate'' means to publish or republish a Document, in whole or
in part, as part of another Document.

An MMC is ``eligible for relicensing'' if it is licensed under this
License, and if all works that were first published under this License
somewhere other than this MMC, and subsequently incorporated in whole
or in part into the MMC, (1) had no cover texts or invariant sections,
and (2) were thus incorporated prior to November 1, 2008.

The operator of an MMC Site may republish an MMC contained in the site
under CC-BY-SA on the same site at any time before August 1, 2009,
provided the MMC is eligible for relicensing.

\subsubsection{ADDENDUM: How to use this License for your documents}

To use this License in a document you have written, include a copy of
the License in the document and put the following copyright and
license notices just after the title page:

\bigskip
\begin{quote}
    Copyright \copyright{}  YEAR  YOUR NAME.\\
    Permission is granted to copy, distribute and/or modify this document
    under the terms of the GNU Free Documentation License, Version 1.3
    or any later version published by the Free Software Foundation;
    with no Invariant Sections, no Front-Cover Texts, and no Back-Cover Texts.
    A copy of the license is included in the section entitled ``GNU
    Free Documentation License''.
\end{quote}
\bigskip

If you have Invariant Sections, Front-Cover Texts and Back-Cover Texts,
replace the ``with \dots\ Texts.''\ line with this:

\bigskip
\begin{quote}
    with the Invariant Sections being LIST THEIR TITLES, with the
    Front-Cover Texts being LIST, and with the Back-Cover Texts being LIST.
\end{quote}
\bigskip

If you have Invariant Sections without Cover Texts, or some other
combination of the three, merge those two alternatives to suit the
situation.

If your document contains nontrivial examples of program code, we
recommend releasing these examples in parallel under your choice of
free software license, such as the GNU General Public License,
to permit their use in free software.


\end {document}
